\cotefile{53}
\scanpage{16}
\[
\left\{\begin{aligned}
&\text{Groupes sur des bases } S\ \text{quelconques. Prendre base } S\ \deleted{\unsure{?}}\ \text{ou prendre base } k\ \text{?}\\
&\text{Groupes \unsure{non-affines|n.-affines} O.K.}
\end{aligned}\right.
\]

Représentations linéaires des $D(M)$ [cas particuliers : $\unsure{\mathbb{G}_m|G_m}$, $\mu_n$]

Application : homomorphismes $D(M)\longrightarrow\mathbb{G}_m$ : on trouve que $M_S$ est ``réflexif'' vis à vis de $\mathbb{G}_m$.

\uline{Dualité de Cartier}\quad $\underline{\mathrm{Hom}}(G,\mathbb{G}_m)$ \deleted{\unsure{$\overset{\Delta}{=}$|est}} représentable ? Cela l'est si $G$ diagonalisable, mais \deleted{pas}\inserted{en général (\unsure{=|ex.|i.e.} pour $G=\mathbb{G}_m$)} pour un schéma en groupes affine (\unsure{même} de t.f.). Par contre, si $G$ est fini loc. libre sur $S$, on a
% slanted
\reversemarginpar\margin{$\underline{\mathrm{Hom}}_{\mathrm{gr}}(G,H)$ représentable si $G$ est fini loc. libre, $H$ affine (disons)}\normalmarginpar
\[
\underline{\mathrm{Hom}}_{\mathrm{gr}}(G,\mathbb{G}_m)
\subset \mathrm{Hom}_{\unsure{\mathcal O|O|\mathbb O}\text{-alg}}
(\mathcal O_S[T,T^{-1}],\unsure{\mathcal A|A})\simeq W(\mathcal A)^*,
\]
et le sous foncteur est celui des \deleted{$\unsure{g}$} $u\in W(\mathcal A)^*(S')$ tels que
\[
\unsure{\Pi|\Delta}(u)=u\otimes u\qquad \varepsilon(u)=1
\]
\margin{[\,\uline{NB} On peut donc oublier la condition $u$ inv. car $\underline{\mathrm{Hom}}_{\mathrm{gr}}(G,\mathbb{G}_m)=\underline{\mathrm{Hom}}_{\unsure{\mathrm{mon}}}(G,\mathcal O^{\unsure{\times|*}})$.}

\deleted{\unsure{La}} La représentabilité est claire, mais de plus on reconnaît que c'est \deleted{\unsure{le|à}} $\mathrm{Spec}(\check{\mathcal A})$, $\check{\mathcal A}$ \deleted{étant}\inserted{$=\underline{\mathrm{Dist}}(G)_{\unsure{/S|S}}$} une algèbre grâce à la structure de coalgèbre de $\mathcal A$. D'autre part, la structure de coalgèbre de $\check{\mathcal A}$ (provenant du fait que le foncteur qu'il représente est un foncteur-groupe
% slanted
\reversemarginpar\margin{$\check{\mathcal A}$ comm. car $G$ l'est. Si $G$ n'est pas comm., il faut prendre le foncteur $k'\longmapsto\mathrm{Hom}_{\unsure{k\text{-alg}}}(\check{\mathcal A},k')=\mathrm{Hom}_{k\text{-alg}}({\check{\mathcal A}\deleted{\unsure{)}}}_{\unsure{\mathrm{ab}}},k')$}\normalmarginpar

\uline{Th.} Tout groupe comm. fini loc. libre \inserted{$G$} est réflexif pour
\par\noindent\begin{tabular}{@{}|p{.91\linewidth}@{}}
$\mathbb{G}_m$. Son dual $G^*$ est fini loc. libre comm., d'où une \uline{\unsure{d'involution|involution}} $G\unsure{\mathrel{\vert\!-\!{\cdot}}|\mapsto}G^*$ de la catégorie des groupes comm. loc. libres. Elle correspond à la dualité des bialgèbres bicommutatives sur $S$.
\end{tabular}

\[
\begin{aligned}
\text{\uline{Ex}}\qquad &\mu_n\longleftrightarrow\mathbb Z/n\mathbb Z
&\ \Bigg|\, &\text{Mais aussi }\ D(M)\longleftrightarrow M_S\\
&\alpha_p\longleftrightarrow\alpha_p && 
\end{aligned}
\]

Marche $\pm$ pour les schémas en groupes \deleted{\unsure{affines}} \deleted{\illegible{1}} sur $S$ \deleted{\unsure{tel que|tels que}\ldots}\inserted{qui} se dévissent en groupes finis loc. libres, groupes discrets, groupes diagonalisables, et versions \deleted{\unsure{discre|discrè}} tordues\ldots\ Si on veut inclure le groupe $\mathbb{G}_a$, il faut \ldots

La dualité de Cartier domine la théorie des groupes algébriques \unsure{(ou} formels) commutatifs, affines ou non -- mais il faut la complémenter par \unsure{considération|la considération} des $\underline{\mathrm{Ext}}^1(G,\mathbb{G}_m)$ alors.
% archive sheet number: 15
\clearpage
